\chapter{Conclusion}\label{ch:conclusion}

%% ===================================================================
\section{Practical Implications}\label{sec:implications}
%% ===================================================================

The findings of this thesis have direct consequences for several practitioner communities.

\paragraph{Model selection and deployment.}
In multi-agent systems where one model's output feeds another's input, a confident but incorrect intermediate result cascades into downstream failures.
The A-R-I profiles provide actionable selection criteria: admittance for uncertainty-critical pipeline stages, resistance for adversarial-facing contexts, and inductance for reasoning-intensive queries.
The rank flow analysis (Figure~\ref{fig:rank-flow}) shows that a model's description quality ranking does not predict its comprehension or behavioural ranking, reinforcing that model selection must be task-aware.

\paragraph{Accessibility tool development.}
The numerical precision bottleneck (mean 2.7 incorrect values per figure) means VLM-generated descriptions cannot yet be trusted without verification for tasks requiring exact values.
Accessibility tools should present descriptions with appropriate confidence caveats, particularly for pie charts and when serving multilingual communities where quality variation is substantial.

\paragraph{Benchmark design.}
Single-score benchmarks systematically miss behavioural dimensions.
Model rankings shuffle across evaluation axes, meaning a leaderboard based on description quality alone misrepresents deployment suitability.
Future benchmarks should incorporate behavioural probes alongside accuracy metrics, and cross-lingual claims require controlled parallel figure designs.

%% ===================================================================
\section{Limitations}\label{sec:limitations}
%% ===================================================================

\paragraph{Sample coverage.}
Human evaluation covers four models on 30 English figures (159 annotations, 39 double-annotated, ICC\,=\,.91); broader coverage across all 13 models and four languages would strengthen judge validation.
The controlled cross-lingual comparison uses only 13 parallel figures, and only 9 inferable elements exist for passive inductance, limiting the granularity of both analyses.

\paragraph{Scoring artefacts.}
Three systematic artefacts affect scoring: colour synonym false positives (${\sim}5\%$ English penalty), dense atoms inflating pie chart penalties by ${\sim}11$ points, and a verbosity penalty that may disadvantage models producing the thorough descriptions accessibility users need.
GPT-4o assigns 78\% Major severity vs Mistral's 47\% on identical errors; while rankings are robust ($\rho_{\text{sys}}\geq.95$), absolute scores depend on judge selection and per-figure scores carry substantial uncertainty ($\rho{\approx}0.6$).

\paragraph{Computational cost.}
Initial experiments evaluated all 1\,005 figures across 13 models, producing ${\sim}$13\,000 base evaluations, but API costs escalated to approximately \$2\,500 for baselines alone before completing a single research question.
Iterative methodology changes (prompt revisions, judge prompt updates, scoring formula adjustments) compounded costs further, as each change required partial re-evaluation.
We therefore adopted a stratified sample (120 atom figures for RQ1, 45 adversarial figures for RQ2--4) with targeted ablations to maintain generalisability within a feasible budget.

%% ===================================================================
\section{Ethical Considerations}\label{sec:ethics}
%% ===================================================================

\paragraph{Data and bias.}
All 1\,005 figures are drawn from publicly available scientific papers on arXiv and open-access repositories; no personally identifiable information appears in the dataset.
Ground-truth annotations were produced by trained annotators following documented guidelines with informed consent.
VLM performance varies substantially across languages, with Bulgarian and Chinese figures receiving lower-quality descriptions, which has direct implications for equitable access to scientific content.
We report these disparities transparently to inform deployment decisions.

\paragraph{Evaluation and deployment.}
The dual-judge evaluation inherits biases from GPT-4o and Mistral Large 3, including the severity split documented in Chapter~\ref{ch:discussion}.
We mitigate this through judge averaging and human validation, but acknowledge that the evaluation framework reflects the perspectives of two specific commercial models.
Even the best-performing model produces ${\sim}2.7$ incorrect numerical values per figure; we caution against deploying VLM-generated descriptions as authoritative without human verification, particularly in contexts where numerical precision is critical.

%% ===================================================================
\section{Future Work}\label{sec:future-work}
%% ===================================================================

\paragraph{Dataset and evaluation extensions.}
Splitting dense atoms into single-fact units would eliminate the pie chart granularity artefact.
Extending human annotation to all 13 models across four languages would enable per-language judge validation.
Expanding the controlled cross-lingual set beyond 13 figures would enable per-chart-type cross-lingual analysis.

\paragraph{Dynamic evaluation with retrieval.}
Future work could evaluate retrieval-augmented generation, where models access the source paper alongside the figure, testing whether additional context improves accuracy or introduces further bias.

\paragraph{Reasoning vs parametric recall.}
GPT-5.2 correctly fabricates 9 of 35 non-inferable blurred elements whose values should not be deducible from context alone.
Controlled experiments with synthetic figures absent from training corpora would disentangle genuine contextual reasoning from parametric recall.

\paragraph{Temporal tracking and user studies.}
Longitudinal evaluation across model versions would reveal whether precision and behavioural reliability improve with successive releases.
Complementary user studies with visually impaired researchers would reveal whether MQM-optimised descriptions actually serve accessibility needs.

%% ===================================================================
\section{Summary of Contributions}\label{sec:contributions}
%% ===================================================================

This thesis makes five contributions to the evaluation of vision-language models on scientific figures.

\begin{enumerate}
\item \textbf{The SciFig benchmark.}
  A multilingual evaluation corpus of 1\,005 scientific figures across four typologically diverse languages (English, Bulgarian, Chinese, German), drawn from native-language venues, with 2\,252 atomic fact checklists, 630 capability questions, and a 45-figure adversarial subset with nine visual transforms.

\item \textbf{Atomic MQM.}
  A normalised scoring framework adapting MQM to figure description, with atom-level tracking, severity-capped weighting, and dual-judge evaluation validated against human annotators (ICC\,=\,.91, $\rho_{\text{sys}}\geq.95$).

\item \textbf{The A-R-I behavioural framework.}
  A three-axis model (Admittance, Resistance, Inductance) that distinguishes faithful readers from silent fabricators and reveals profiles invisible to accuracy metrics.

\item \textbf{A 13-model empirical evaluation.}
  The largest comparative evaluation of VLMs on scientific figure understanding to date, spanning 4B to 235B+ parameters with 26\,364 errors across 2\,966 evaluations.

\item \textbf{Cross-cutting findings.}
  Numerical precision is the dominant bottleneck (30.6\% of errors). Description, comprehension, and behaviour are partially dissociated capabilities. Robustness is independent of clean-condition accuracy. Per-language comparisons require controlled parallel designs.
\end{enumerate}

\bigskip
\noindent
Returning to the four research questions posed in Chapter~\ref{chap:introduction}:
\textbf{RQ1} finds that GPT-5.2 leads (75.5 MQM) but the top four models cluster within 5 points; German's advantage is a dataset artefact.
\textbf{RQ2} shows rotation and low contrast cause the largest degradation; robustness is distinct from baseline accuracy.
\textbf{RQ3} reveals comprehension and description engage different capacities, with Gemini overtaking GPT-5.2 on question answering.
\textbf{RQ4} demonstrates that A-R-I exposes behavioural patterns invisible to accuracy metrics: GPT-5.2's resistance gap, Gemma's silent fabrication, and the active-vs-passive inductance inversion.

Together, these findings argue that evaluating VLMs on scientific figures requires moving beyond single-score benchmarks toward multi-dimensional assessment of accuracy, robustness, comprehension, and behavioural trustworthiness.
